\documentclass[12pt]{article}
\usepackage[margin=1in]{geometry}
\usepackage{amsmath, amssymb}
\usepackage{graphicx}
\usepackage{booktabs}
\usepackage{hyperref}
\usepackage{parskip}
\usepackage{titlesec}

\title{\textbf{NoisyNetworks: Exploration via Learned Noise}\\
\large Asterix Atari Environment --- DQN, Double DQN, and PER}
\author{}
\date{}

\begin{document}
\maketitle
\tableofcontents
\newpage

% ============================================================
\section{Motivation}
% ============================================================

A core challenge in reinforcement learning is the \textbf{exploration-exploitation tradeoff}.
The agent must try new actions to discover better policies, but also exploit what it already knows.

The dominant approach in deep RL --- $\varepsilon$-greedy --- is simple: with probability
$\varepsilon$ take a random action, otherwise act greedily.
Its problems:
\begin{itemize}
  \item Exploration is \emph{undirected}: random actions ignore the current knowledge of the value function.
  \item $\varepsilon$ must be annealed by hand; the schedule is a hyperparameter that interacts with everything else.
  \item All actions are explored equally regardless of state, wasting samples in well-understood regions.
\end{itemize}

\textbf{NoisyNets} (Fortunato et al., 2017) replace the fixed $\varepsilon$-greedy policy with
\emph{learned, state-dependent noise} injected directly into the network weights,
letting the network itself decide how much to explore in each situation.

% ============================================================
\section{Background: DQN}
% ============================================================

\subsection{Q-Learning and the Bellman Equation}

Q-learning seeks to approximate the optimal action-value function:
\begin{equation}
  Q^*(s, a) = \mathbb{E}\!\left[r + \gamma \max_{a'} Q^*(s', a') \;\Big|\; s, a\right]
\end{equation}

DQN (Mnih et al., 2015) parameterises $Q(s,a;\theta)$ with a deep neural network and
minimises the temporal-difference (TD) loss:
\begin{equation}
  \mathcal{L}(\theta) =
  \mathbb{E}_{(s,a,r,s') \sim \mathcal{D}}\!\left[
    \Bigl(
      \underbrace{r + \gamma \max_{a'} Q(s', a';\,\theta^-)}_{\text{TD target}}
      - Q(s,a;\,\theta)
    \Bigr)^2
  \right]
\end{equation}

where $\theta^-$ are the \textbf{target network} parameters (periodically copied from $\theta$)
and $\mathcal{D}$ is the \textbf{replay buffer}.

\subsection{Key Stabilisation Tricks}
\begin{itemize}
  \item \textbf{Experience replay}: store $(s,a,r,s')$ tuples; sample mini-batches uniformly.
        Breaks temporal correlations between consecutive updates.
  \item \textbf{Target network}: a frozen copy of $Q$ used only for computing TD targets.
        Prevents the moving-target problem.
\end{itemize}

% ============================================================
\section{Background: Double DQN}
% ============================================================

\subsection{The Overestimation Problem}

Standard DQN uses the \emph{same} network to both \emph{select} and \emph{evaluate} the
greedy action in the target:
\[
  y^{\text{DQN}} = r + \gamma \max_{a'} Q(s', a';\,\theta^-)
\]
Because $\max$ is applied to noisy estimates, this systematically \emph{overestimates}
$Q^*$. The agent pursues inflated Q-values rather than truly better states.

\subsection{Double DQN Correction}

Double DQN (van Hasselt et al., 2015) decouples action \emph{selection} (online network)
from action \emph{evaluation} (target network):
\begin{equation}
  y^{\text{DDQN}} = r + \gamma\, Q\!\left(s',\; \arg\max_{a'} Q(s',a';\,\theta);\; \theta^-\right)
\end{equation}

The online network picks the best action; the target network scores it.
Overestimation is reduced because both networks must agree for a value to be inflated.

% ============================================================
\section{Background: Prioritised Experience Replay (PER)}
% ============================================================

Uniform replay wastes capacity: the agent replays ``easy'' transitions (low TD error)
as often as ``hard'' ones. PER (Schaul et al., 2015) samples proportional to TD error.

\subsection{Priority and Sampling Probability}

The priority of transition $i$ is:
\begin{equation}
  p_i = |\delta_i| + \epsilon_{\text{PER}}
\end{equation}
where $\delta_i$ is the absolute TD error and $\epsilon_{\text{PER}}$ is a small constant
preventing zero probability.

The sampling probability is:
\begin{equation}
  P(i) = \frac{p_i^\alpha}{\sum_k p_k^\alpha}
\end{equation}

$\alpha \in [0,1]$ controls how strongly priorities are used ($\alpha = 0$ recovers uniform).

\subsection{Importance Sampling Correction}

Prioritised sampling introduces bias because the distribution no longer matches the
on-policy distribution. Importance sampling (IS) weights correct this:
\begin{equation}
  w_i = \left(\frac{1}{N \cdot P(i)}\right)^{\!\beta}
  \Big/ \max_j w_j
\end{equation}

$\beta$ is annealed from a small value to 1 over training. The TD update becomes:
\begin{equation}
  \mathcal{L}(\theta) = \mathbb{E}\!\left[w_i \cdot \delta_i^2\right]
\end{equation}

% ============================================================
\section{Background: NoisyNets}
% ============================================================

\subsection{Noisy Linear Layers}

A standard linear layer computes $y = Wx + b$.
A \textbf{noisy linear layer} replaces each weight and bias with a learned mean plus
learned-variance noise:
\begin{equation}
  y = \bigl(\mu^w + \sigma^w \odot \varepsilon^w\bigr)\,x
    + \bigl(\mu^b + \sigma^b \odot \varepsilon^b\bigr)
\end{equation}

where $\varepsilon^w, \varepsilon^b$ are random noise tensors sampled at each forward pass,
and $\mu, \sigma$ are trainable parameters.

\subsection{Factorised Gaussian Noise}

To reduce the number of independent noise samples (expensive for large layers),
\textbf{factorised} noise uses two independent vectors $\varepsilon_i^p$ (input) and
$\varepsilon_j^q$ (output) and constructs the weight noise as:
\begin{equation}
  \varepsilon^w_{ij} = f(\varepsilon_i^p)\, f(\varepsilon_j^q),
  \qquad
  f(x) = \operatorname{sgn}(x)\sqrt{|x|}
\end{equation}

This reduces the noise parameters from $O(pq)$ to $O(p + q)$.

\subsection{Parameter Initialisation}

\[
  \mu^w_{ij} \sim \mathcal{U}\!\left[-\tfrac{1}{\sqrt{p}},\, \tfrac{1}{\sqrt{p}}\right],
  \qquad
  \sigma^w_{ij} = \frac{\sigma_0}{\sqrt{p}}
\]
where $p$ is the number of inputs and $\sigma_0 = 0.5$ by default.

\subsection{Why This Beats $\varepsilon$-Greedy}

\begin{itemize}
  \item Noise is \textbf{state-dependent}: the same weights are active across a full episode,
        so exploration is \emph{consistent} (the agent commits to exploratory behaviour rather
        than flipping randomly at every step).
  \item The network \textbf{learns to reduce noise} when it is confident and maintain noise
        when it is uncertain, via the gradient of the loss on $\sigma$.
  \item No schedule to tune: $\varepsilon$ is eliminated entirely.
\end{itemize}

% ============================================================
\section{Experimental Setup}
% ============================================================

\begin{itemize}
  \item \textbf{Environment}: Asterix (Atari 2600 via ALE), frame-stacked, grayscale.
  \item \textbf{Variants tested}: Vanilla DQN, Vanilla DDQN, Vanilla DQN+PER,
        NoisyNet-DQN, NoisyNet-DDQN, NoisyNet-DQN+PER.
  \item \textbf{Training}: 20M environment steps per variant.
  \item \textbf{Metric}: mean episodic reward over the last 100 episodes.
  \item \textbf{Network}: convolutional backbone (3 conv layers) $\rightarrow$ 2 noisy
        (or standard) linear layers.
\end{itemize}

% ============================================================
\section{Results and Analysis}
% ============================================================

\subsection{NoisyNet-DQN vs.\ Vanilla DQN}

\textbf{Finding}: The two curves are nearly identical for the entire 20M step run.
Both converge to $\approx 27$--$29$ mean reward.

\textbf{Why}: In a dense-reward environment like Asterix, $\varepsilon$-greedy with a
simple linear annealing schedule already provides sufficient coverage of the state space.
The learned noise in NoisyNet essentially replicates what $\varepsilon$-greedy was already
doing --- it has no underexplored regions to find that the baseline missed.
NoisyNets shine most in environments where exploration \emph{direction} matters;
here, random actions are good enough.

\subsection{NoisyNet-DDQN vs.\ Vanilla DDQN}

\textbf{Finding}: NoisyNet-DDQN clearly outperforms Vanilla DDQN after $\sim$5M steps,
converging to $\approx 30$ vs.\ $\approx 27$.

\textbf{Why}: DDQN reduces Q-value overestimation. With more accurate value estimates,
the noisy exploration signal is \emph{meaningful} --- the network explores regions where
it is genuinely uncertain rather than chasing inflated Q-values.
In Vanilla DQN, overestimation corrupts the signal that NoisyNet would use to guide
exploration; once DDQN corrects that bias, NoisyNet's learned uncertainty maps cleanly
onto actual policy improvements.

\subsection{Vanilla DDQN vs.\ Vanilla DQN (Baseline)}

\textbf{Finding}: Vanilla DDQN does \emph{not} clearly outperform Vanilla DQN.
Both converge to $\approx 25$--$28$.

\textbf{Why}: DDQN's correction for overestimation matters most in environments where
inflated Q-values cause the agent to commit to \emph{systematically wrong} actions.
Asterix has a dense, consistent reward structure: even overestimated Q-values tend to point
toward genuinely good actions (collecting items is reliably rewarded).
The overestimation bias exists but does not distort the \emph{ranking} of good vs.\ bad
actions enough to degrade the final policy.
In harder, sparser environments (e.g., MuJoCo locomotion, sparse navigation) DDQN's
advantage would likely be more pronounced.

\subsection{NoisyNet-PER vs.\ Vanilla PER}

\textbf{Finding}: NoisyNet-PER outperforms Vanilla PER by a consistent margin ($\approx 28$--$29$
vs.\ $\approx 24$--$25$).

\textbf{Why}: PER and NoisyNets are \emph{synergistic}:
\begin{enumerate}
  \item NoisyNet exploration generates diverse, surprising transitions --- these have
        high TD errors and get prioritised by PER.
  \item PER then replays these informative transitions more often, amplifying the
        learning signal from NoisyNet's targeted exploration.
  \item The combination creates a positive feedback loop: better exploration $\rightarrow$
        richer replay buffer $\rightarrow$ more efficient learning.
\end{enumerate}
Without NoisyNet, Vanilla PER's replay buffer eventually fills with stale transitions
whose TD errors have already been driven down. There is no mechanism to keep generating
fresh high-priority experience.

\subsection{Vanilla PER: Fast Start, Early Plateau}

\textbf{Finding}: Vanilla PER learns the fastest in the 0--5M step window but plateaus
sooner and slightly lower than Vanilla DQN in the long run.

\textbf{Why}: PER's prioritised sampling is extremely efficient early --- high-error
transitions are valuable and get heavily replayed. However, as the agent improves,
TD errors shrink uniformly and the ``hard'' transitions become indistinguishable from easy ones.
The replay buffer ages: frequently replayed transitions become stale as the network
moves away from the policy that generated them.
Without a mechanism (like NoisyNet) to continuously generate new high-priority experiences,
PER runs out of steam and ε-greedy's uniform exploration keeps feeding DQN with fresher data.

\subsection{Summary Table}

\begin{center}
\begin{tabular}{lcc}
  \toprule
  Variant & Approx.\ Final Reward & Key Observation \\
  \midrule
  Vanilla DQN         & 27--29 & Strong baseline \\
  NoisyNet-DQN        & 27--29 & No gain over vanilla \\
  Vanilla DDQN        & 25--28 & No clear gain over DQN \\
  NoisyNet-DDQN       & $\approx$ 30 & Best late-game noisy variant \\
  Vanilla PER         & 24--25 & Fast early, plateaus low \\
  NoisyNet-PER        & 28--29 & Fast + sustained improvement \\
  \bottomrule
\end{tabular}
\end{center}

% ============================================================
\section{Key Takeaways}
% ============================================================

\begin{enumerate}
  \item \textbf{NoisyNets are not universally better than $\varepsilon$-greedy.}
        In dense-reward environments with simple action spaces, the directed exploration
        provides no measurable benefit over random perturbation.

  \item \textbf{NoisyNets require accurate value estimates to be effective.}
        Paired with DDQN (lower overestimation), NoisyNet's uncertainty signal maps to
        real policy gaps. Paired with standard DQN, the overestimated Q-landscape
        dilutes the benefit.

  \item \textbf{PER and NoisyNets are complementary.}
        PER makes replay efficient; NoisyNet keeps the replay buffer populated with
        fresh, high-priority transitions. Neither alone achieves what both together do.

  \item \textbf{Algorithmic improvements compound.}
        The biggest gains come from stacking orthogonal techniques (DDQN addresses
        overestimation; PER addresses sample efficiency; NoisyNet addresses exploration).

  \item \textbf{Environment matters.}
        Asterix is relatively forgiving (dense rewards, clear feedback). Results may
        differ significantly on sparse-reward or procedurally generated environments
        (e.g., MiniGrid) where directed exploration is the binding constraint.
\end{enumerate}

% ============================================================
\section{Equations Reference (Quick Copy)}
% ============================================================

\[
  \text{DQN target:} \quad
  y^{\text{DQN}} = r + \gamma \max_{a'} Q(s', a';\,\theta^-)
\]

\[
  \text{DDQN target:} \quad
  y^{\text{DDQN}} = r + \gamma\, Q\!\left(s',\; \arg\max_{a'} Q(s',a';\,\theta);\; \theta^-\right)
\]

\[
  \text{PER priority:} \quad p_i = |\delta_i| + \varepsilon, \quad
  P(i) = \frac{p_i^\alpha}{\sum_k p_k^\alpha}
\]

\[
  \text{IS weight:} \quad w_i = \left(\frac{1}{N \cdot P(i)}\right)^{\!\beta} \Big/ \max_j w_j
\]

\[
  \text{Noisy layer:} \quad
  y = (\mu^w + \sigma^w \odot \varepsilon^w)\,x + (\mu^b + \sigma^b \odot \varepsilon^b)
\]

\[
  \text{Factorised noise:} \quad
  \varepsilon^w_{ij} = f(\varepsilon_i)\,f(\varepsilon_j), \quad f(x) = \operatorname{sgn}(x)\sqrt{|x|}
\]

\end{document}
